% ---
% title: GUMBEL-SOFTMAX
% date: 2025/12/18
% tag: ML
% ---
\definecolor{mdc1}{HTML}{FFFFFF}
\definecolor{mdc2}{HTML}{FF0000}
This post is a summary of \href{https://arxiv.org/pdf/1611.00712.pdf}{Gumbel-Softmax}, which solves the issue of categorical sampling that is not differentiable.

\section{数学原理：离散采样怎么变成可反传？}

\subsection{问题：Categorical 采样不可导}

你有一个分类分布（K 类）：

\begin{itemize}
  \item logits：\(a \in \mathbb{R}^K\)
  \item 概率：\(\large \pi_i = \frac{e^{a_i}}{\sum_j e^{a_j}}\)
\end{itemize}

真正的采样是：

\[
y \sim \text{Categorical}(\pi), \quad y \in \{e_1,\dots,e_K\}
\]

这是离散的 one-hot，\textbf{无法对 logits 做梯度反传}。

\subsection{Gumbel-Max Trick：离散采样的等价形式}

先定义 Gumbel 噪声：

\[
g_i \sim \text{Gumbel}(0,1),\quad g_i = -\log(-\log u_i),\ u_i\sim \text{Uniform}(0,1)
\]

则有经典结论（Gumbel-Max trick）：

\[
\arg\max_i (a_i + g_i)\ \sim\ \text{Categorical}(\text{softmax}(a))
\]

也就是说，\colorbox{mdc2}{\textcolor{mdc1}{给每个 logit 加 Gumbel 噪声再 argmax，就等价于从 categorical 采样。}}

但是 argmax 还是不可导。

\subsection{Gumbel-Softmax：把 argmax 换成 softmax（可导近似）}

把不可导的 argmax 用 softmax 近似，引入温度 \(\tau>0\)：

\[
\tilde{y}_i = \frac{\exp\left((a_i + g_i)/\tau\right)}{\sum_j \exp\left((a_j + g_j)/\tau\right)}
\]

则 \(\tilde{y}\) 是一个 \textbf{“软 one-hot”}（分量非负且和为 1），并且对 \(a\) \textbf{可导}。

关键性质：

\begin{itemize}
  \item \(\tau \to 0\)：\(\tilde{y}\) 趋近 one-hot（更“离散”，但梯度更不稳定）
  \item \(\tau \to \infty\)：\(\tilde{y}\) 趋近均匀分布（更“平滑”，但偏差大）
\end{itemize}

这就是用 \textbf{重参数化（reparameterization）} 把采样写成：

\[
\tilde{y}=f(a, u),\quad u\sim U(0,1)
\]

从而可以反传到 \(a\)。

\subsection{Straight-Through (ST) Gumbel-Softmax（训练更像真离散）}

很多时候你希望：

\begin{itemize}
  \item 前向：真的 one-hot（像离散决策）
  \item 反向：用软的 \(\tilde{y}\) 给梯度
\end{itemize}

做法：

\[
y^{hard}=\text{onehot}(\arg\max \tilde{y})
\]

反传时用 trick：

\[
y = y^{hard} - \text{stopgrad}(\tilde{y}) + \tilde{y}
\]

等价于：\textbf{前向用 hard，梯度当作 soft}。

\section{伪代码流程（标准版 + ST 版）}

\subsection{标准 Gumbel-Softmax（soft 输出）}

\begin{lstlisting}[language=python]
Input: logits a (K-dim), temperature tau
Sample u_i ~ Uniform(0,1) for i=1..K
Compute g_i = -log(-log(u_i))        # Gumbel noise
Compute z_i = (a_i + g_i) / tau
Output y_soft = softmax(z)          # differentiable "soft one-hot"
\end{lstlisting}

\subsection{Straight-Through Gumbel-Softmax（hard 前向，soft 梯度）}

\begin{lstlisting}[language=python]
Input: logits a, temperature tau
y_soft = GumbelSoftmax(a, tau)
k = argmax(y_soft)
y_hard = onehot(k)

# forward uses y_hard, backward uses y_soft
Output y = y_hard - stopgrad(y_soft) + y_soft
\end{lstlisting}

\section{简单例子：三选一“离散开关”如何可导训练？}

假设你要从 3 个操作里选 1 个（比如三种滤波器 / 三个专家 / 三个 token）：

\begin{itemize}
  \item logits \(a=[2.0,\ 0.0,\ -1.0]\)
  \item softmax 概率大概偏向第 1 类
\end{itemize}

\subsection{如果直接采样}

你会得到 one-hot，比如 [1,0,0]，但因为采样/argmax 不可导，loss 无法把梯度传回 logits。

\subsection{用 Gumbel-Softmax（\textbf{\(\tau=0.5\)}）}

采样一次 \(u\) 得到 gumbel 噪声 g，然后：\(\tilde{y}=\text{softmax}((a+g)/0.5)\)

可能得到类似：\(\tilde{y}=[0.92,\ 0.07,\ 0.01]\)

这不是严格 one-hot，但已经很“尖锐”，并且\textbf{对} a \textbf{可导}。

\subsection{用 ST 版本}

\begin{itemize}
  \item 前向用 [1,0,0]（真的选了第 1 类）
  \item 反向梯度按 [0.92,0.07,0.01] 传播
  
  你就能像训练连续网络那样训练“离散选择”。
\end{itemize}
